\documentclass[UTF8,fontset=fandol,aspectratio=169,11pt]{ctexbeamer}
% Shared Beamer preamble for Big Data and Management Decision-Making slides.

\usetheme{Madrid}
\usecolortheme{default}
\usefonttheme{professionalfonts}

\usepackage{amsmath,amssymb,mathtools}
\usepackage{booktabs,array}
\usepackage{tikz}
\usepackage{graphicx}
\usepackage{listings}
\usepackage{xcolor}
\usetikzlibrary{arrows.meta,positioning,shapes.geometric}

\newcommand{\BDMFandolPath}{/usr/local/texlive/2025/texmf-dist/fonts/opentype/public/fandol/}
\setCJKmainfont[
  Path=\BDMFandolPath,
  UprightFont=FandolSong-Regular.otf,
  BoldFont=FandolSong-Regular.otf,
  ItalicFont=FandolKai-Regular.otf,
  BoldItalicFont=FandolKai-Regular.otf
]{FandolSong-Regular.otf}
\setCJKsansfont[
  Path=\BDMFandolPath,
  UprightFont=FandolSong-Regular.otf,
  BoldFont=FandolSong-Regular.otf
]{FandolSong-Regular.otf}
\setCJKmonofont[
  Path=\BDMFandolPath,
  UprightFont=FandolSong-Regular.otf,
  BoldFont=FandolSong-Regular.otf
]{FandolSong-Regular.otf}
\setmonofont{Menlo}

\definecolor{AcademicBlue}{RGB}{22,44,73}
\definecolor{AcademicTeal}{RGB}{22,119,119}
\definecolor{AcademicGray}{RGB}{76,84,95}
\definecolor{AcademicLight}{RGB}{245,247,250}
\definecolor{AcademicAmber}{RGB}{181,111,35}
\definecolor{CodeBg}{RGB}{248,248,248}

\setbeamercolor{structure}{fg=AcademicBlue}
\setbeamercolor{frametitle}{fg=white,bg=AcademicBlue}
\setbeamercolor{title}{fg=white,bg=AcademicBlue}
\setbeamercolor{block title}{fg=white,bg=AcademicBlue}
\setbeamercolor{block body}{bg=AcademicLight,fg=black}
\setbeamertemplate{navigation symbols}{}
\setbeamertemplate{footline}[frame number]
\setbeamertemplate{itemize item}{\raisebox{0.2ex}{\(\blacktriangleright\)}}
\setbeamerfont{title}{series=\mdseries}
\setbeamerfont{subtitle}{series=\mdseries}
\setbeamerfont{frametitle}{series=\mdseries}
\setbeamerfont{block title}{series=\mdseries}
\setbeamerfont{section in toc}{series=\mdseries}

\lstdefinestyle{pythonstyle}{
  basicstyle=\ttfamily\scriptsize,
  backgroundcolor=\color{CodeBg},
  keywordstyle=\color{AcademicBlue},
  commentstyle=\color{AcademicGray},
  stringstyle=\color{AcademicAmber},
  showstringspaces=false,
  breaklines=true,
  frame=single,
  rulecolor=\color{black!20},
  columns=fullflexible
}

\lstdefinestyle{sqlstyle}{
  language=SQL,
  basicstyle=\ttfamily\scriptsize,
  backgroundcolor=\color{CodeBg},
  keywordstyle=\color{AcademicBlue},
  commentstyle=\color{AcademicGray},
  stringstyle=\color{AcademicAmber},
  showstringspaces=false,
  breaklines=true,
  frame=single,
  rulecolor=\color{black!20},
  columns=fullflexible
}

\lstset{style=pythonstyle}

\hypersetup{
  colorlinks=true,
  linkcolor=AcademicBlue,
  urlcolor=AcademicTeal,
  pdfauthor={大数据与管理决策基础},
  pdfsubject={大数据与管理决策基础课程课件}
}


\title[企业数据与业务对象]{第2讲：企业数据与业务对象}
\subtitle{关系模型、ER 模型、星型模型与对象映射}
\author{《大数据与管理决策基础》}
\date{}

\begin{document}

\begin{frame}
  \titlepage
  \begin{center}
    \small 核心命题：企业数据建模不是为了画表，而是为了让指标、模型和行动具有稳定对象基础。
  \end{center}
\end{frame}

\begin{frame}{本讲学习目标}
\begin{enumerate}
  \item 区分实体、属性、关系、事件、状态与指标；
  \item 设计订单--客户--产品--库存 ER 模型；
  \item 理解事实表粒度、主键、外键和完整性约束；
  \item 将关系表组织为星型模型；
  \item 将表映射为业务对象、关系和可执行行动。
\end{enumerate}
\end{frame}

\begin{frame}{企业运营数据的五类基本概念}
\begin{table}
\small
\centering
\begin{tabular}{p{0.16\linewidth}p{0.34\linewidth}p{0.36\linewidth}}
\toprule
类型 & 含义 & 示例\\
\midrule
实体 & 具有稳定身份的对象 & 客户、产品、供应商、仓库\\
属性 & 实体或事件的描述变量 & 地区、类别、单价、客户等级\\
事件 & 时间上发生的一次业务活动 & 下单、发货、付款、退货\\
状态 & 某一时点对象的状态快照 & 库存量、订单状态、设备状态\\
指标 & 对事件或状态的聚合度量 & 销售额、准时交付率\\
\bottomrule
\end{tabular}
\end{table}
\end{frame}

\begin{frame}{关系模型：表是有限关系}
若客户集合为 \(C\)，产品集合为 \(P\)，时间集合为 \(T\)，则订单数据可表示为
\[
O\subseteq C\times P\times \mathbb N\times T.
\]

\begin{block}{三个基本约束}
\begin{enumerate}
  \item 实体完整性：主键不能为空且唯一；
  \item 引用完整性：外键必须引用已存在记录；
  \item 业务完整性：数量、价格、日期满足业务规则。
\end{enumerate}
\end{block}
\end{frame}

\begin{frame}{订单--客户--产品--库存 ER 模型}
\begin{center}
\resizebox{0.86\linewidth}{!}{%
\begin{tikzpicture}[
  entity/.style={draw=AcademicBlue!60, fill=AcademicLight, rounded corners=1pt,
                 minimum width=2.7cm, minimum height=0.9cm, align=center},
  arrow/.style={-{Stealth[length=2mm]}, thick, AcademicTeal}
]
\node[entity] (cust) {Customer\\客户};
\node[entity, right=3cm of cust] (order) {Order\\订单};
\node[entity, right=3cm of order] (prod) {Product\\产品};
\node[entity, below=1.5cm of prod] (inv) {Inventory\\库存快照};
\draw[arrow] (cust) -- node[above]{1:N} (order);
\draw[arrow] (prod) -- node[above]{1:N} (order);
\draw[arrow] (prod) -- node[right]{1:N} (inv);
\end{tikzpicture}
}
\end{center}
\begin{alertblock}{关键问题}
一行订单事实到底表示“订单头”、还是“订单明细项”？粒度决定后续 KPI 是否可解释。
\end{alertblock}
\end{frame}

\begin{frame}{星型模型：面向分析的组织方式}
\[
\text{FactOrders}
\leftarrow
\{\text{DimCustomer},\text{DimProduct},\text{DimDate},\text{DimChannel}\}.
\]

\begin{table}
\small
\centering
\begin{tabular}{p{0.22\linewidth}p{0.58\linewidth}}
\toprule
组件 & 设计要点\\
\midrule
事实表 & 明确粒度，保存度量值和维度外键。\\
维度表 & 描述分析切片，如客户、产品、日期、渠道。\\
指标口径 & 基于事实表粒度定义销售额、毛利率、延迟率。\\
\bottomrule
\end{tabular}
\end{table}
\end{frame}

\begin{frame}{事实表粒度与度量可加性}
\begin{table}
\small
\centering
\begin{tabular}{p{0.23\linewidth}p{0.34\linewidth}p{0.29\linewidth}}
\toprule
概念 & 判断问题 & 示例\\
\midrule
粒度 & 一行记录代表什么业务事实？ & 一个订单明细项。\\
可加度量 & 能否沿所有维度求和？ & 销售金额、数量。\\
半可加度量 & 是否不能跨时间直接相加？ & 库存余额。\\
不可加度量 & 是否需要由分子分母重算？ & 准时率、毛利率。\\
\bottomrule
\end{tabular}
\end{table}

\begin{alertblock}{建模规范}
先声明粒度，再选择维度和事实；不同粒度的事实不要混入同一张事实表。
\end{alertblock}
\end{frame}

\begin{frame}{事件与状态不能混淆}
\begin{columns}[T]
\begin{column}{0.48\linewidth}
\begin{block}{事件表}
\begin{itemize}
  \item 下单、发货、付款；
  \item 有明确发生时间；
  \item 适合计算流量、频次、周期。
\end{itemize}
\end{block}
\end{column}
\begin{column}{0.48\linewidth}
\begin{block}{状态快照表}
\begin{itemize}
  \item 库存余额、订单状态；
  \item 描述某一时点状态；
  \item 适合计算存量、风险暴露。
\end{itemize}
\end{block}
\end{column}
\end{columns}
\end{frame}

\begin{frame}{从关系表映射为业务对象}
\begin{table}
\small
\centering
\begin{tabular}{p{0.24\linewidth}p{0.35\linewidth}p{0.28\linewidth}}
\toprule
业务对象 & 来源表 & 可支持行动\\
\midrule
Customer & customers & 分层运营、优惠券、回访\\
Product & products & 定价、促销、补货\\
Order & orders + dimensions & 加急、拆单、客户沟通\\
InventoryPosition & inventory + products & 补货、调拨、清仓\\
\bottomrule
\end{tabular}
\end{table}

\[
\text{Tables}\rightarrow\text{Objects}\rightarrow\text{Links}\rightarrow\text{Actions}.
\]
\end{frame}

\begin{frame}[fragile]{配套代码：生成 DDL 与数据字典}
\begin{lstlisting}[language=Python]
@dataclass(frozen=True)
class Entity:
    name: str
    primary_key: str
    fields: list[Field]
    description: str

def create_table_sql(entity: Entity) -> str:
    ...
\end{lstlisting}

\begin{block}{运行方式}
\begin{lstlisting}[language=bash]
PYTHONPATH=src python3 src/bdm_decision/cases/week02_enterprise_data_model.py
\end{lstlisting}
\end{block}
\end{frame}

\begin{frame}{课堂练习}
以“零售订单--客户--产品--库存”为背景：
\begin{enumerate}
  \item 列出至少四个实体及其主键；
  \item 为每个实体设计至少四个字段；
  \item 标明实体之间的基数关系；
  \item 指出一个事实表和至少三个维度表；
  \item 将至少两个关系表映射为业务对象，并说明可执行行动。
\end{enumerate}
\end{frame}

\begin{frame}{常见错误}
\begin{itemize}
  \item 不区分订单头与订单明细，导致事实表粒度混乱；
  \item 使用业务含义不稳定的字段作为主键；
  \item 只画 ER 图，不说明主外键和完整性约束；
  \item 将库存快照误认为交易事件；
  \item 在 KPI 计算前没有定义指标口径和事实表粒度。
\end{itemize}
\end{frame}

\end{document}
